\section{Gestión del replay buffer y selección de datos}

\subsection{Por qué son clave el tamaño y la estrategia de muestreo}

El replay buffer proporciona los datos para múltiples rondas de entrenamiento, pero su diseño influye directamente en el desempeño. Las dos dimensiones centrales son la \textbf{capacidad} y la \textbf{estrategia de muestreo}:

\begin{itemize}
    \item \textbf{Capacidad}: si es demasiado pequeña, sobrescribe con frecuencia las muestras antiguas y pierde diversidad; si es demasiado grande, reduce la eficiencia del entrenamiento. En la práctica, suele elegirse una capacidad que pueda almacenar varios millones de transitions.
    \item \textbf{Estrategia de muestreo}: el muestreo uniforme es el más simple; Prioritized Experience Replay permite que las transitions con TD error alto se visiten con más frecuencia.
\end{itemize}

\begin{knowledgebox}{Efecto directo de la estrategia de muestreo en el desempeño}
Las transitions con TD error alto suelen incluir decisiones más difíciles o escenarios de colapso; priorizar su retorno permite que la función Q converja más rápido, pero si la priorización es excesiva, puede introducir un sesgo en la estimación. El compromiso habitual es usar pesos de importance-sampling con un $\beta$ que decae hacia 1 para corregir el sesgo.
\end{knowledgebox}

\subsection{Limpieza de los datos del replay: filtrado frente a complementación}

Cada vez que el agente ejecuta una nueva política, debe evaluarse la calidad de los datos recién recopilados:

\begin{itemize}
    \item Filtrado: retirar las transitions con reward anómalo o fuera de los límites (por ejemplo, obtener un reward de 1 tras chocar con una pared), para evitar ruido de entrenamiento;
    \item Complementación: sintetizar artificialmente transitions con etiquetas específicas (por ejemplo, reward shaping, domain randomization) para ayudar al agente a aprender situaciones raras pero cruciales.
\end{itemize}

\begin{warningbox}{Efecto de contaminación del replay buffer}
Si el buffer acumula datos de baja calidad durante un periodo prolongado (por ejemplo, las exit transitions tras el colapso de la policy), el modelo puede seguir optimizándose sobre ejemplos incorrectos. Vaciar periódicamente cerca del 7\% de las muestras con el reward más bajo del buffer y volver a recolectarlas puede evitar el «data poisoning».
\end{warningbox}

\subsection{Resumen de la sección}

La capacidad del replay buffer, la estrategia de muestreo y la estrategia de limpieza de datos determinan en conjunto la estabilidad del aprendizaje off-policy. Un muestreo prioritario moderado, la corrección IS y la eliminación/el muestreo de las transitions anómalas pueden evitar que el modelo sobreajuste sobre datos antiguos o se vea afectado por el ruido.

